\section{Summary and conclusions}\label{Sec:Conc}

In this article, we have presented the production of strange and non-strange light-flavor particles at midrapidity (\etarange 0.8) in high-multiplicity pp collisions at \s $=$ 13 TeV, as a function of the unweighted transverse spherocity \SOPT. In order to account for the biases between weakly-decaying and charged primary hadrons, we have elaborated on how the observable was updated from the traditional transverse spherocity.

The multiplicity was estimated in two different kinematic regions: at midrapidity (\etarange 0.8) by  measuring the activity in the SPD (\tracklet), and at forward rapidities ($2.8 < \eta < 5.1$ and $-3.7 < \eta < -1.7$) by measuring the activity in both sides of the V0 forward detector (V0M amplitude). In Fig.~\ref{Fig:CL1vsV0M}, we report that the different \SOPT  event classes have a large variance in local charged particle density when estimating the multiplicity through the V0M, while being roughly equal in terms of \mpt. In contrast, the \tracklet-selected \SOPT events are constrained in terms of $\langle \rm{d}\it{N}_{\pi}/\rm{d}\it{y} \rangle$ within $\sim 10\%$, while having a large variance in \mpt. We conclude that estimating multiplicity at midrapidity allows one to isolate and study the dynamics of particle production in events that are driven by either soft or hard-QCD physics.

The main features of this analysis are reported in the particle-to-pion ratios in Figs.~\ref{Fig:CombRatioPy}--~\ref{Fig:CombRatioHerEx}, highlighting an enhancement of strange hadrons in events with an isotropic topology, and a strong suppression in events with a jet-like topology. This indicates that events with an isotropic topology describe the average high-multiplicity event fairly well, while events heavily influenced by jet-like physics are outliers. For the baryon-to-meson ratios presented in Fig.~\ref{Fig:Flow}, the \SOPT selection also highlights an enhancement at intermediate \pt in the p/\PI, \LA/\kzero and \XI/\PHI ratios, but without a depletion of low-\pt particles for jet-like events. The effects are consistent among the three ratios, and the \PHI exhibits a behavior similar to the other strange particles. However, the \pt-integrated yield of \PHI as a function of \SOPT presented in Fig.~\ref{Fig:CL1_10_Int_phi} highlights features similar to other non-strange particles, reporting no observed modification of relative \PHI meson production as a function of \SOPT.

The presented models are not able to describe the quantitative observations found in data for the particle ratios, but are overall able to describe the interplay between \SOPT classified events at high-multiplicity with good accuracy in the double ratios. Remarkably, even though the production mechanisms for the PYTHIA 8.2 Monash and PYTHIA 8.2 Ropes variations (ropes formed by layers of overlapping strings) are qualitatively different, they both predict the interplay with only minor differences between them. 

Finally, we report the relative integrated strange particle yield to pions as function of \SOPT in Fig.~\ref{Fig:CL1_1_Int_py}. It is found that one can achieve a similar strangeness enhancement found in multiplicity differential analyses, by fixing the local charged particle density and varying the azimuthal topology, instead of only varying the charged-particle density. This indicates that particle production at high multiplicities is driven by more than a single source, with different strangeness-to-pion production rates. The same measurement is then performed in high-multiplicity events selected at forward rapidities. Fig.~\ref{Fig:CL1vsV0M} highlights that the differential-\SOPT selection in this event-class primarily varies the local charged particle density approximately a factor of two. Therefore, when contrasted to the multiplicity-dependent ALICE measurements in Ref.~\cite{Naturepaper}, one could naively expect to observe strangeness suppression and enhancement due to the difference in charged particle density alone. Surprisingly, we show in Fig.~\ref{Fig:V0M_1_Int} that this effect is negligible when estimating multiplicity in the forward-rapidity region, with no significant strangeness enhancement or suppression within systematic uncertainties. This goes contrary to expectations based on previous ALICE publications, and the effect is currently not well understood. PYTHIA 8.2 Monash and Herwig 7.2 are unable to describe this effect, but PYTHIA 8.2 Ropes and EPOS-LHC qualitatively capture the trend, showcasing a large suppression of strangeness production in jet-like events, although with a mass ordering that is incompatible with the experimental data. 

We conclude that with \SOPT, one is able to categorize events into classes based on characteristic azimuthal topologies, from jet-like events associated with hard-QCD physics, to isotropic events associated with soft-QCD physics. Furthermore, the findings presented in this article suggest that average high-multiplicity events, even in the most extreme cases, are dominated by soft processes, where rare hard processes play little or no role for bulk observables.

